\documentclass[10pt]{article}
\usepackage[margin=0.95in]{geometry}
\usepackage{amsmath,amssymb,amsthm,microtype,xcolor}
\usepackage[colorlinks=true,urlcolor=blue!50!black]{hyperref}
\newtheorem{theorem}{Theorem}
\setlength{\emergencystretch}{3em}
\title{Classification of Positive Natural Collatz Seeds\\with Mechanical Parity Itineraries}
\author{Collatz proof project}
\date{September 5, 2026}
\begin{document}
\maketitle
\begin{abstract}
We complete the repository's mechanical-word branch. For arbitrary real
slope and intercept, a positive natural seed with a globally mechanical
shortcut Collatz parity itinerary must be 1 or 2, and its slope must be
$1/2$. Every positive orbit with a mechanical tail therefore reaches 1.
The proof combines a factor-complexity bound, exact rational phase coding,
and the endpoint-swap cycle obstruction. All real intercepts and
nonprimitive rational presentations are handled. No universal assertion
that Collatz itineraries become mechanical is proved. This is a
formalization and integration of classical ingredients, not a proof of
Collatz or a claim of mathematical novelty.
\end{abstract}

\section{Statements}
Let $T(n)=n/2$ for even $n$ and $(3n+1)/2$ for odd $n$. The predicate
$\operatorname{HasItin}(\alpha,\rho,N)$ means that, at every natural time $t$,
\[
 T^t(N)\bmod2
 =\lfloor(t+1)\alpha+\rho\rfloor-\lfloor t\alpha+\rho\rfloor.
\]
\begin{theorem}
For all $\alpha,\rho\in\mathbb R$ and $N\in\mathbb N$ with $N>0$,
$\operatorname{HasItin}(\alpha,\rho,N)$ implies
$\alpha=1/2$ and $N\in\{1,2\}$.
\end{theorem}
\begin{theorem}
If $N>0$ and $\operatorname{HasItin}(\alpha,\rho,T^s(N))$ for some
natural $s$ and real $\alpha,\rho$, then $T^t(N)=1$ for some natural $t$.
\end{theorem}
The positivity condition is essential: zero has the all-zero itinerary.
The first theorem restricts seeds and slope; it does not assert that every
intercept realizes both seeds.

\section{Periods and rotations are actual returns}
If $q\alpha$ is an integer, translating the two floors by this integer
shows that the mechanical word is $q$-periodic. Naturals with identical
infinite parity itineraries are equal, by congruence modulo $2^L$ for
every $L$. Hence $T^q(N)=N$ for any natural seed with this itinerary.
In particular, a rational denominator is a return period from the initial
seed, without a boundedness assumption.

A realized return word rotated by $s$ is realized by $T^s(N)$.
This is proved by reducing indices modulo the period and pumping the
return. The endpoint-swap theorem therefore applies to rotations: if a
return word has both $1u0$ and $0u1$ as rotations, then $u$ is empty and
the corresponding shifted seeds are 1 and 2.

\section{A uniform finite rational grid}
For $0<p<q$ and $\gcd(p,q)=1$, define the carry bit
\[
 b_r(t)=\mathbf1\{(tp+r)\bmod q+p\ge q\}.
\]
Multiplication by $p$ permutes residues modulo $q$. Thus, from any phase
$r$, some shifts reach phase 0 and phase $q-1$.
For $0<i<q$, the residue $ip\bmod q$ is nonzero. This gives the following
structure for the two extreme length-$q$ words:
\[
 (b_{q-1}(0),\ldots,b_{q-1}(q-1))=1u0,\qquad
 (b_0(0),\ldots,b_0(q-1))=0u1.
\]
Indeed, the first letters are 1 and 0 and the last letters are 0 and 1.
At interior positions, replacing the residue by one less cannot change
the carry: equality at the threshold would make $(i+1)p$ divisible by $q$.
The grid is periodic, so both shifted words are actual return words.
Endpoint cancellation now forces a visit to 1.

\section{All real intercepts and all rational presentations}
Normalize $\rho$ to $\beta=\{\rho\}$ and put $r=\lfloor q\beta\rfloor$.
The inequalities $r\le q\beta<r+1$, together with quotient and remainder
arithmetic, prove for every natural $t$ that
\[
 \left\lfloor t\frac pq+\beta\right\rfloor
 =\left\lfloor\frac{tp+r}{q}\right\rfloor.
\]
Successive differences are exactly $b_r(t)$. This is an equality of
infinite words, including negative original intercepts and threshold
boundary cases, not a finite or floating-point approximation.
Consequently every reduced rational mechanical itinerary in $(0,1)$
reaches 1.

For an arbitrary real slope, the earlier mechanical complexity bound
$p_N(L)\le L+1$ forces boundedness. Repeated states give a positive return
period $q$ on a tail. Mechanical telescoping over $k$ repeats gives
$-1<k(q\alpha-p)<1$, where $p$ is the odd-step count. Hence $q\alpha=p$.
Positive-cycle inequalities give $0<p<q$. Dividing $p,q$ by their gcd
produces a reduced rational slope, so the grid theorem applies to this tail
and the original seed reaches 1. No rational presentation is assumed
primitive in this step.

Finally shift to the visit to 1. Its two-step return and telescoping force
$2\alpha=1$. This integral slope increment gives a two-step return at the
original seed. The checked two-step-return classification gives $N=1$ or 2.

\section{Formal artifact and validation}
The new modules are \texttt{MechanicalPeriod.lean},
\texttt{RationalGrid.lean}, and \texttt{RationalMechanical.lean} under
\texttt{lean/Collatz/}. The principal declarations are
\texttt{mechanical\_reaches\_one}, \texttt{mechanical\_classification}, and
\texttt{eventual\_mechanical\_reaches\_one} in namespace \texttt{Collatz}.
The modular phase, endpoint structure, real-intercept transport, and period
transport are separately proved and reusable.

The Lean build and selected axiom audit pass: 107 audited declarations
use only the standard foundations. Four exact computational controls pass,
including negative and boundary intercepts and the two permitted seeds.
These tests are regression checks; the universal classification is a Lean
proof. The audit is dependency inspection, not independent kernel replay.

\section{Context and the remaining Collatz gap}
The complexity ingredient is related to Dubickas,
\emph{On integer sequences generated by linear maps} (2009), Theorem 5,
\href{https://doi.org/10.1017/S0017089508004655}{doi:10.1017/S0017089508004655}.
The cycle cancellation follows Knight's high-cycle argument,
\emph{Collatz high cycles do not exist} (2026),
\href{https://doi.org/10.1016/j.disc.2025.114812}{doi:10.1016/j.disc.2025.114812},
also inspected in the nested Knight formalization. No nested module or
custom mathematical axiom is imported.

The previously documented rational rotation and intercept bridges are now
closed. General Collatz counterexamples need not have mechanical tails.
Proving that every positive orbit eventually has such a tail would itself
settle the conjecture and cannot be assumed as a coverage shortcut.
General escape and nontrivial cycles remain unresolved outside this class.
\end{document}
